
%\documentclass[./main.tex(path of main file)]{subfiles}
\documentclass[10pt]{article}
\usepackage{amsmath}
\DeclareMathOperator*{\argmin}{arg\,min} % thin space, limits underneath in displays
\DeclareMathOperator*{\argmax}{arg\,max} % thin space, limits underneath in displays
\newtheorem{thm}{Theorem}
\usepackage{amssymb}
\usepackage{amsfonts}
\usepackage{mathrsfs}
\usepackage{bm}
\usepackage{indentfirst}
\setlength{\parindent}{0em}
\usepackage[margin=1in]{geometry}
\usepackage{graphicx}
\usepackage{setspace}
\doublespacing
\usepackage[flushleft]{threeparttable}
\usepackage{booktabs,caption}
\usepackage{float}
\usepackage[sort,comma]{natbib}
\usepackage[hidelinks]{hyperref}
\usepackage{booktabs}
\usepackage{multirow}

\usepackage{import}
\usepackage{xifthen}
\usepackage{pdfpages}
\usepackage{transparent}
\usepackage{subfiles}
\usepackage{blindtext}

\newcommand{\incfig}[1]{%
\def\svgwidth{\columnwidth}
\import{./figures/}{#1.pdf_tex}
}




\title{}
\author{}
\date{}


\begin{document}
\maketitle


\section{Inference about partial effect}
When there is a conflict between the t-ratio for coef of variable x
and one for average partial effect (APE) of variable x, trust the
one for the former (t-ratio for coef of x), not APE.
The implication runs in the direction from the structure (coef) to
the partial effects, not the reverse.
The use of the standard error for the partial effect is, perhaps,
not for hypothesis test, but for developing confidence intervals.
This is introduced in Green's book on page 754 (Example 17.9 
Hypothesis Tests About Partial Effects).

Even though we prefer to use Average Partial Effect (APE), however, 
APE is slightly different from Partial Effect at the mean (PEA).
Hence, it is okay to use PEA when the dataset is larger, as it
requires large computational power.

Note, the single estimate of the APE might not capture the 
interesting variation in the partial effect as other variables
change. We need to plot the APE for variable x as it varies
with another variable y. Add the confidence interval for 
the APE of x. For example, I can plot the APE for LFE component 
(tau) as it varies with inflation expectation (PX1).



Note, the marginal effect is no longer for a 1 unit change even though most people would
interpret it that way when using marginal effects. It is the derivative!!! It reflects
the change in probability given a very small change in the variable of interest (In stata, the
small change is 0.0001)
However, you can manually compute the marginal effect given a change of variable is one.

\section{Inference about odds ratios }

\begin{align*}
\text{ Odds ratio } &= exp(coef)\\
\text{ Standard Error of Odds ratio } &= \text{ Odds ratio }  \times \text{ SE of coef }\\
\text{ t-ratio of Odds ratio } &= \text{ t-ratio of coef }\\
\text{ p-value of Odds ratio } &= \text{ p-value of coef }\\
\text{ Confidence interval (CI) of Odds ratio } &= exp(\text{ CI of coef })
\end{align*}


\section{Interaction effects}
The partial effect of an interaction reported is useless.(Page 756 Green's book)
Consider regression below
\begin{equation*}
Prob(DocVis_{it} > 0) = F(b_1 + b_2Age_{it} + b_3Edu_{it} + b_4Age_{it} \times Edu_{it}).
\end{equation*}
The partial effect of $ Age $ should be $ F'(\bm{x}'\bm{b})(b_2 + b_4Edu) $,
where $ F() $ and $ F'() $ are cdf and pdf of a probability distribution.
We need to manually compute the partial effect following the method mentioned above, 
instead of using the one reported in common computer packages.




\section{Pooled estimator in a panel data}

The pooled MLE that ignores fixed effects will be inconsistent--possibly wildly so.
(Note: Because the estimator is ML, not least squares, converting the data to deviations
from group means is not a solution--converting the binary dependent variable to deviations
will produce a new variable with unknown properties.)(page 781 in Green's book)







\section{Multinomial choices}

Conditional logit model (Multinomial logit model) is different from ordered logit model.

\subsection{Conditional logit model}
McFadden has shown that if (and only if) the $ J $ disturbances are independent and
identically distributed with Gumbel distribution,
\begin{equation*}
F(\varepsilon_{ij}) = exp( - exp( - \varepsilon_{ij})),
\end{equation*}
then
\begin{equation*}
Prob(Y_{i} = j) = \frac{exp(\bm{z'}_{ij} \bm{\theta})}{\sum\limits_{j = 1} ^J	
exp(\bm{z'}_{ij} \bm{\theta})},
\end{equation*}
which leads to what is called the conditional logit model. (It is often labeled the
multinomial logit model, but this wording conflicts with the usual name for the model
discussed in the next section.)

{\textbf {The coefficients in this model are difficult to interpret.}}
For any particular exogenous variable x, the partial effect need not have the same sign
as the coef of it. SE can be estimated using the delta method.



\bibliographystyle{plainnat}
\bibliography{my_bib}

\end{document}

